% CIE 6004 Assignment 2 report
\documentclass[10pt,twocolumn,letterpaper]{article}
\usepackage[review]{cvpr}
\input{preamble}
\definecolor{cvprblue}{rgb}{0.21,0.49,0.74}
\usepackage[pagebackref,breaklinks,colorlinks,allcolors=cvprblue]{hyperref}
\def\paperID{*****}
\def\confName{CVPR}
\def\confYear{2026}

\title{Reconstructing Garden, Reconstructing an Object,\\and Placing It Back: An Empirical 3DGS Study}
\author{Student Name\\
Student ID\\
\tt\small student@example.com}

\begin{document}
\maketitle

\begin{abstract}
We build a complete 3D Gaussian Splatting (3DGS) pipeline for the three tasks of Assignment~2:
reconstructing the Garden scene, reconstructing one supplied object (an Olympus OM-2 camera)
from handheld clips, and placing that object back into Garden with a consistent depth ordering.
On the 12 public validation views our Garden model reaches \textbf{27.27\,dB PSNR / 0.8668 SSIM /
0.0764 LPIPS}, improving on the released baseline ($25.73 / 0.7925 / 0.1757$) on all three
metrics. For the object we build a disjoint 285/33/164 frame split, run a mask-limited SfM
that registers all 498 frames, and train a white-background 3DGS model. A controlled experiment
on the mask boundary ($+2$\,px dilation, everything else fixed) raises masked PSNR from
16.23 to 18.23\,dB while leaving one rear view unsolved, which localises the remaining error to
view coverage rather than to mask bias alone. Finally we place the camera on the garden table,
using a RANSAC tabletop plane and a width-based scale estimate, and render a 10\,s orbit whose
depth ordering is produced by the splatting itself rather than by compositing.
\end{abstract}

\section{Setup and conventions}

\paragraph{Data.} Garden~\cite{barron2022} ships 161 training photographs at $1297\times840$ with a COLMAP
reconstruction whose stored camera model is $5187\times3361$ ($f_x\!\approx\!3844.9$).
Because focal length and resolution scale together, the field of view is unchanged if the images
are loaded at their native size, and all projections must use $f_x\!\approx\!961.4$ at
$1297\times840$. Poses are \texttt{world\_to\_camera} $4\times4$ matrices in OpenCV convention
with column vectors. We verified the convention by projecting the 135\,521 sparse points into a
validation photograph before training; the table rim and door align to sub-pixel error
.

\paragraph{Models.} Garden and the object are trained with the official 3DGS implementation~\cite{kerbl2023}
(commit \texttt{54c035f}) with $30$k iterations and SH degree 3, on one H800 PCIe (80\,GB),
20 vCPU and 100\,GB RAM. Garden runs in 21 minutes at full resolution with 4.17\,M Gaussians;
the object trains in 5 minutes on 285 frames of $406\times720$.

\paragraph{Camera handling.} Our renderer builds \texttt{Camera(R, T, FoVx, FoVy)} directly from
the JSON poses with $R=M_{3\times3}^{\top}$ (camera-to-world) and $T=M_{3\times3+1}$, matching the
convention of the COLMAP loader. Getting this transposition wrong yields a plausible-looking but
mis-oriented render (12.6\,dB PSNR), which is why every stage below is gated by a numeric check.

\section{Task A: Garden}

\begin{table}[t]
\centering\small
\setlength{\tabcolsep}{4pt}
\begin{tabular}{lccc}
\hline
 & PSNR$\uparrow$ & SSIM$\uparrow$ & LPIPS$\downarrow$\\
\hline
Released baseline & 25.73 & 0.7925 & 0.1757\\
Ours (30k, full res) & \textbf{27.27} & \textbf{0.8668} & \textbf{0.0764}\\
\hline
\end{tabular}
\caption{Average over the 12 validation views, measured with the official
\texttt{tools/evaluate.py} (LPIPS-Alex v0.1~\cite{zhang2018}).}
\label{tab:garden}
\end{table}

We train on all 161 views with automatic holdout disabled, at native resolution, and render the
12 validation cameras and the 12 test cameras independently of the training loader. The
prescribed video is rendered from the 240 poses of \texttt{trajectory/garden.json} at
$1280\times720$ and 24\,fps, i.e.\ exactly 10.0\,s; these cameras have a different intrinsics
($f=1410$) from the evaluation cameras, so they are constructed separately.

Table~\ref{tab:garden} reports the validation metrics. All three improve over the baseline;
the largest relative gain is in LPIPS ($-57\%$), consistent with sharper foliage and wood grain.

\paragraph{The hardest view.} DSC07988 is a steep, downward-looking view of the tabletop
(Fig.~\ref{fig:hardest}). It reaches 20.44\,dB against a per-view mean of 27.27, i.e.\ a 6.8\,dB
gap. Two effects compound: the camera sees the plank seams almost edge-on, where a small depth
error shifts a high-frequency stripe by several pixels, and the reconstruction has relatively few
training views from this elevation. The 2$\times$ crops show that our model keeps the radial seam
pattern but smooths the thin flower stems, while the baseline loses the seam structure entirely.

\begin{figure}[t]
\centering
\includegraphics[width=\linewidth]{fig/hardest_view.jpg}
\caption{Hardest validation view (DSC07988): reference, ours, released baseline.}
\label{fig:hardest}
\end{figure}

\begin{figure}[t]
\centering
\includegraphics[width=\linewidth]{fig/hardest_crop.jpg}
\caption{2$\times$ crops of the vase and stems. Our seam structure survives; thin stems are
smoothed. The baseline blurs the seam pattern.}
\label{fig:crop}
\end{figure}

\section{Task B: the camera}

\subsection{Split and poses}
The supplied clips have 198, 136 and 164 frames at 8\,fps. We use 1975 and 1976 for development
and reserve 1977 for testing, and cut a contiguous validation arc from 1976 with a one-second
(8-frame) gap on both sides: train 285 frames (all of 1975, plus 1976 \#0--31 and \#81--135),
validation 33 frames (1976 \#40--72), test 164 frames (all of 1977). The split is written to
\texttt{object\_split.json} with explicit frame identifiers and a \texttt{val\_gap\_frames} field.

Automatic SfM on the full set with plain SIFT slows down quadratically and, more importantly,
reconstructs the plaid bed rather than the object. We therefore extract features \emph{inside the
foreground masks} (COLMAP \texttt{ImageReader.mask\_path}), which restricts the reconstruction to
the camera body. All 498 frames register in one model with 10\,939 points and a final
reprojection cost of 0.61\,px. Masks come from \texttt{u2netp}; 12 random frames were inspected
before use.

\subsection{Training and the defect}
The object is trained on white-composited crops of $406\times720$ with 3DGS (30k iterations,
131\,972 Gaussians). The baseline orbit (Fig.~\ref{fig:triplet}) shows one clear defect: the
silhouette is \emph{eaten away} along the dark underside and the right-hand body edge, and a
faint white halo floats beside the body. Our explanation is mask-boundary bias: the hard mask
treats the low-contrast boundary pixels as background, so the photometric residual never asks the
model to cover them, and the freed capacity drifts into the background as floaters.

\begin{figure}[t]
\centering
\includegraphics[width=\linewidth]{fig/mask_triplet.jpg}
\caption{Object pipeline: source frame, foreground mask, novel-view render (baseline).}
\label{fig:triplet}
\end{figure}

\paragraph{Controlled experiment.} We test exactly one variable: dilating every training mask by
2\,px before compositing, keeping the split, the initial point cloud, the seed and the 30k budget
fixed. Masked PSNR over three held-out validation views rises from 16.23 to 18.23\,dB, and the
front view improves from 20.63 to 23.90\,dB. Silhouette IoU against the \emph{undilated}
reference mask drops slightly (0.834 to 0.811), which is expected: the variant deliberately covers
two extra pixels, so this number is biased against the change and we report it rather than hide it.

\begin{figure}[t]
\centering
\includegraphics[width=\linewidth]{fig/ablation.jpg}
\caption{Baseline (left) versus $+2$\,px mask dilation (right) at three fixed validation views.}
\label{fig:ablation}
\end{figure}

\paragraph{What did not improve.} One rear/underside view stays at about 11\,dB in both variants.
Dilation cannot help there: the failure follows the viewpoints, not the mask. The capture orbits
the object but rarely looks at the bottom face, so those Gaussians are supported by a handful of
grazing observations. This is the remaining failure we report.

\begin{figure}[t]
\centering
\includegraphics[width=\linewidth]{fig/stages.jpg}
\caption{The same validation camera at 5k, 10k, 15k and 30k iterations.}
\label{fig:stages}
\end{figure}

\begin{figure}[t]
\centering
\includegraphics[width=\linewidth]{fig/loss.png}
\caption{Training loss (total and L1), 30k iterations; the kink near 15k is the end of
densification.}
\label{fig:loss}
\end{figure}

\section{Task C: placing the camera in Garden}

\paragraph{Scene geometry.} We fit the tabletop directly in the reconstruction: points whose
projection falls inside an image-space ellipse around the table are selected in a validation view,
and RANSAC returns a plane with 8\,409 inliers. Its normal (within $4.1^{\circ}$ of the average
camera up vector over six views) defines the scene up axis, its inlier centroid the table centre,
and the inlier spread a radius of 0.99 scene units; the camera sits 0.78 units above the plane.
The two independent estimates agree, which fixes the metric scale to roughly one unit per metre.

\paragraph{Scale and transform.} The object's pruned point cloud has a PCA frame in which the body
width, height and depth measure 5.28, 2.99 and 2.13 units. An OM-2 body is 136\,mm wide, so the
uniform scale is $s=0.090$ and the resulting object is 13.6\,cm wide, 7.7\,cm tall and 5.5\,cm deep.
The transform is a similarity: positions are mapped by $sR p + t$, log-scales are offset by
$\log s$, and rotations are left-multiplied by the quaternion of $R$; this is exactly what
\texttt{fusion/placement.json} stores as a $4\times4$ matrix acting on column vectors.

\paragraph{Occlusion.} The camera is placed on the table 0.45 units from the vase, with its lens
axis toward the first frame of the shot. The 240-frame, 10\,s path arcs $185^{\circ}$ around the
table at 1.9\,units radius and 1.26\,units height. It starts with the vase overlapping the object,
which makes the first second a genuine depth-ordering test, and the arc then slides the vase off
the line of sight until the camera stands fully visible against the garden. Depth ordering is not
composited: the object
Gaussians are merged into the scene point cloud, so the rasteriser resolves which splat is in
front at every pixel.

\begin{figure}[t]
\centering
\includegraphics[width=\linewidth]{fig/fusion_strip.jpg}
\caption{Fusion frames at 0.0\,s, 5.0\,s and 10.0\,s: the camera rests on the table, the vase
overlaps it at the start of the arc, and it is fully visible at the end.}
\label{fig:fusion}
\end{figure}

\section{Bonus: making the insertion look placed}

The basic render shows three artefacts: no contact shadow, a brightness mismatch between the
object and the table, and background Gaussians bleeding through the object's silhouette. We
improve the first two: (i) table Gaussians within a radius of the object base have their SH DC
term darkened by a quadratic falloff, which produces a soft contact shadow without new geometry;
(ii) the object's SH DC is rescaled toward the local table illuminant measured in the same frame.
The before/after pair is rendered from the same camera and placement.

\section{Reproducibility and limitations}

All runs are scripted (\texttt{code/}), with fixed seeds, recorded configs
(\texttt{runs/*/cfg\_args}) and per-run logs. The repository commit, CUDA and PyTorch versions,
hardware and wall-clock times are listed in \texttt{README.md}. The main limitation is the object
capture itself: the rear and bottom faces are poorly observed, so no amount of mask tuning
recovers them, and the halo that survives pruning still shows as a faint streak at some angles.

\bibliographystyle{ieeenat_fullname}
\bibliography{main}

\end{document}
